\par\Needspace{6\baselineskip}\medskip
\hypertarget{native-137-article099}{}\pdfbookmark[2]{99.}{article-099}
\begin{center}99.\end{center}
For the HYPERBOLA we likewise retain the symbols $p$, $v$, $v'$, $f$, $F$, $r$, $r'$, and $t$ with the same meanings. For the semimajor axis $a$, which is negative here, we write $-\alpha$. As in Art. 21 and the following articles, we set the eccentricity $e$ $=\dfrac1{\cos\psi}\mathpunct{\cdot}$ We set the auxiliary quantity there denoted by $u$ $=\dfrac Cc$ for the first position and $=Cc$ for the second. It follows readily that $c$ is always greater than 1, \SourceLeaf{138}but, other things being equal, differs less from 1 the closer the two given positions are to each other. We take the sixth and seventh equations developed in Art. 21 and put them in a slightly different form here.
\[
\begin{aligned}
\mathrm{1)}\qquad&\cos\tfrac12v=\tfrac12\left(\BareRoot{\dfrac Cc}+\BareRoot{\dfrac cC}\right)\BareRoot{\dfrac{(e-1)\alpha}r}\\[.8em]
\mathrm{2)}\qquad&\sin\tfrac12v=\tfrac12\left(\BareRoot{\dfrac Cc}-\BareRoot{\dfrac cC}\right)\BareRoot{\dfrac{(e+1)\alpha}r}\\[.8em]
\mathrm{3)}\qquad&\cos\tfrac12v'=\tfrac12\left(\BareRoot{Cc}+\BareRoot{\dfrac1{Cc}}\right)\BareRoot{\dfrac{(e-1)\alpha}{r'}}\\[.8em]
\mathrm{4)}\qquad&\sin\tfrac12v'=\tfrac12\left(\BareRoot{Cc}-\BareRoot{\dfrac1{Cc}}\right)\BareRoot{\dfrac{(e+1)\alpha}{r'}}\mathpunct{\cdot}
\end{aligned}
\]
\noindent From these we immediately obtain the following equations.
\[
\begin{aligned}
\mathrm{5)}\qquad&\sin F=\tfrac12\alpha\left(C-\dfrac1C\right)\BareRoot{\dfrac{ee-1}{rr'}}\\[.8em]
\mathrm{6)}\qquad&\sin f=\tfrac12\alpha\left(c-\dfrac1c\right)\BareRoot{\dfrac{ee-1}{rr'}}\\[.8em]
\mathrm{7)}\qquad&\cos F=\left(e\left(c+\dfrac1c\right)-\left(C+\dfrac1C\right)\right)\dfrac\alpha{2\BareRoot{rr'}}\\[.8em]
\mathrm{8)}\qquad&\cos f=\left(e\left(C+\dfrac1C\right)-\left(c+\dfrac1c\right)\right)\dfrac\alpha{2\BareRoot{rr'}}\mathpunct{\cdot}
\end{aligned}
\]
\noindent Further, equation X of Art. 21 gives
\[
\begin{aligned}
\frac r\alpha&=\tfrac12e\left(\frac Cc+\frac cC\right)-1\\[.6em]
\frac{r'}\alpha&=\tfrac12e\left(Cc+\frac1{Cc}\right)-1,
\end{aligned}
\]
\noindent and hence
\[
\begin{aligned}
\mathrm{9)}\qquad&\frac{r'-r}\alpha=\tfrac12e\left(C-\frac1C\right)\left(c-\frac1c\right)\\[.6em]
\mathrm{10)}\qquad&\frac{r'+r}\alpha=\tfrac12e\left(C+\frac1C\right)\left(c+\frac1c\right)-2.
\end{aligned}
\]
\noindent Combining this equation 10 with equation 8 gives
\[
\mathrm{11)}\qquad\alpha=\frac{r'+r-\left(c+\dfrac1c\right)\cos f\,\mathpunct{.}\BareRoot{rr'}}{\tfrac12\left(c-\dfrac1c\right)^2}\mathpunct{\cdot}
\]
\noindent As in the ellipse, we therefore set $\dfrac{\BareRoot{\dfrac{r'}r}+\BareRoot{\dfrac r{r'}}}{2\cos f}=1+2l$ or $=1-2L$, \SourceLeaf{139}according as $\cos f$ is positive or negative. This gives
\[
\begin{aligned}
\mathrm{12)}\qquad&\alpha=\frac{8\left(l-\tfrac14\left(\BareRoot c-\BareRoot{\dfrac1c}\right)^2\right)\cos f\,\mathpunct{.}\BareRoot{rr'}}{\left(c-\dfrac1c\right)^2}\\[1em]
\mathrm{12^*)}\qquad&\alpha=\frac{-8\left(L+\tfrac14\left(\BareRoot c-\BareRoot{\dfrac1c}\right)^2\right)\cos f\,\mathpunct{.}\BareRoot{rr'}}{\left(c-\dfrac1c\right)^2}\mathpunct{\cdot}
\end{aligned}
\]
\noindent As in the ellipse, the quantity $l$ or $L$ is calculated with the aid of the auxiliary angle $\omega$. Finally, using hyperbolic logarithms, we obtain from equation XI of Art. 22
\[
\begin{aligned}
\frac{kt}{\alpha^{\frac32}}&=\tfrac12e\left(Cc-\frac1{Cc}-\frac Cc+\frac cC\right)-\log Cc+\log\frac Cc\\[.6em]
&=\tfrac12e\left(C+\frac1C\right)\left(c-\frac1c\right)-2\log c
\end{aligned}
\]
\noindent or, eliminating $C$ with the aid of equation 8,
\[
\frac{kt}{\alpha^{\frac32}}=\frac{\left(c-\dfrac1c\right)\cos f\,\mathpunct{.}\BareRoot{rr'}}\alpha\mathpunct{\raisebox{-1ex}{.}}+\tfrac12\left(cc-\frac1{cc}\right)-2\log c.
\]
\noindent In this equation we substitute for $\alpha$ its value from 12 or $12^*$. We then introduce the symbol $m$ or $M$ with the meaning assigned by formulas 11 and $11^*$ of Art. 88. Finally, for brevity, we write
\[
\begin{gathered}
\tfrac14\left(\BareRoot c-\BareRoot{\frac1c}\right)^2=z,\\[.6em]
\frac{cc-\dfrac1{cc}-4\log c}{\tfrac14\left(c-\dfrac1c\right)^3}=Z,
\end{gathered}
\]
\noindent This gives the equations
\[
\begin{aligned}
\mathrm{13)}\qquad&m=(l-z)^{\frac12}+(l-z)^{\frac32}Z\\[.6em]
\mathrm{13^*)}\qquad&M=-(L+z)^{\frac12}+(L+z)^{\frac32}Z,
\end{aligned}
\]
\noindent which involve only one unknown, $z$, since $Z$ is evidently a function of $z$ given by the following formula.
\[
Z=\frac{(1+2z)\BareRoot{(z+zz)}-\log(\BareRoot{1+z}+\BareRoot z)}{2(z+zz)^{\frac32}}\mathpunct{\cdot}
\]
